\part{Приближенные методы решения}
\section{Глава 9. Предсказание с единой стратегией и аппроксимацией}

\subsection{Упражнение 9.1(с. 252)}
Покажите, что табличные методы типа тех, что представлены в первой части книги, --- частный случай линейной аппроксимации функции. Каковы в этом случае векторы признаков?
\begin{solution}
    Возьмём $d=|\mathcal S|$ и в качестве признаков --- one-hot векторы: для
    состояния $s$
    \begin{equation*}
        x_i(s) = \begin{cases} 1, & i = s,\\ 0, & i \ne s, \end{cases}
        \qquad i = 1, \dots, |\mathcal S|.
    \end{equation*}
    Каждому состоянию отвечает свой базисный вектор, и все они ортонормированы.

    При таких признаках линейная аппроксимация (9.8) вынимает ровно одну
    компоненту весов:
    \begin{equation*}
        \hat v(s, \mathbf w) = \mathbf w^\top \mathbf x(s)
        = \sum_{i=1}^{|\mathcal S|} w_i x_i(s) = w_s.
    \end{equation*}
    То есть вектор весов $\mathbf w$ --- это и есть таблица ценностей, а $w_s$ ---
    ячейка состояния $s$.

    Метод тоже вырождается в табличный. Градиент здесь $\nabla\hat v(s,\mathbf w)
    = \mathbf x(s)$, поэтому общее линейное СГС-обновление
    \begin{equation*}
        \mathbf w_{t+1} = \mathbf w_t + \alpha\big[U_t - \hat v(S_t,\mathbf w_t)\big]\mathbf x(S_t)
    \end{equation*}
    из-за one-hot $\mathbf x(S_t)$ затрагивает единственную компоненту $i=S_t$, а
    все остальные ($i\ne S_t$) оставляет без изменения:
    \begin{equation*}
        w_{S_t} \leftarrow w_{S_t} + \alpha\big[U_t - w_{S_t}\big],
        \qquad w_i \leftarrow w_i \ (i \ne S_t).
    \end{equation*}
    Это дословно табличное обновление с целью $U_t$ (например, $U_t = G_t$ для
    Монте-Карло или $U_t = R_{t+1} + \gamma\hat v(S_{t+1},\mathbf w_t)$ для TD(0)).
    Обновления разных состояний не влияют друг на друга --- ровно потому, что
    признаки ортогональны, --- так что обобщения между состояниями нет, как и
    положено таблице.
\end{solution}

\subsection{Упражнение 9.2(с. 253)}
Почему формула (9.17) определяет $(n+1)^k$ различных признаков для размерности
$k$?
\begin{equation}
    x_i(s) = \prod_{j=1}^{k} s_j^{\,c_{i,j}},
    \tag{9.17}
\end{equation}
\begin{solution}
    Признак задаётся набором показателей $(c_1, \dots, c_k)$, где каждый
    $c_j$ независимо пробегает множество $\{0, 1, \dots, n\}$ из $n+1$ элементов.
    По правилу произведения таких наборов
    \begin{equation*}
        \underbrace{(n+1)(n+1)\cdots(n+1)}_{k} = (n+1)^k.
    \end{equation*}

    Осталось проверить, что разным наборам отвечают разные признаки (иначе
    различных признаков было бы меньше). Пусть два набора задают один и тот же
    моном при всех значениях аргументов:
    \begin{equation*}
        \prod_{j=1}^{k} s_j^{c_j} = \prod_{j=1}^{k} s_j^{c'_j}
        \quad \text{для всех } s_1, \dots, s_k.
    \end{equation*}
    Зафиксируем все аргументы кроме $s_j$, положив их равными $1$. Слева и справа
    остаётся $s_j^{c_j} = s_j^{c'_j}$ при всех $s_j$, откуда $c_j = c'_j$. Это
    верно для каждого $j$, поэтому наборы совпадают.

    Значит, отображение из наборов показателей в признаки инъективно: различных
    наборов $(n+1)^k$, и различных признаков ровно столько же.
\end{solution}

\subsection{Упражнение 9.3 (с. 254)}
При каких $n$ и $c_{i,j}$ порождаются векторы признаков
\begin{equation*}
    \mathbf x(s) = \big(1,\ s_1,\ s_2,\ s_1 s_2,\ s_1^2,\ s_2^2,\ s_1 s_2^2,\ s_1^2 s_2,\ s_1^2 s_2^2\big)^\top ?
\end{equation*}

\begin{solution}
    Здесь $k=2$ (две компоненты $s_1, s_2$), и каждый признак имеет вид
    $x_i(s) = s_1^{c_{i,1}} s_2^{c_{i,2}}$. Признаков ровно $9 = 3^2 = (n+1)^k$,
    причём это все комбинации показателей из $\{0,1,2\}^2$ --- то есть полный
    полиномиальный базис порядка $n=2$ (показатели $c_{i,j} \in \{0,1,2\}$).

    По порядку признаков пары показателей $(c_{i,1}, c_{i,2})$ таковы:
    \begin{equation*}
        \begin{array}{c|ccccccccc}
            x_i(s) & 1 & s_1 & s_2 & s_1 s_2 & s_1^2 & s_2^2 & s_1 s_2^2 & s_1^2 s_2 & s_1^2 s_2^2 \\
            \hline
            (c_{i,1}, c_{i,2}) & (0,0) & (1,0) & (0,1) & (1,1) & (2,0) & (0,2) & (1,2) & (2,1) & (2,2)
        \end{array}
    \end{equation*}
\end{solution}

\subsection{Упражнение 9.4(с. 264)}
Допустим, у нас есть основания полагать, что одно из двух измерений пространства состояний оказывает большее влияние на функцию ценности, чем другое, и что обобщение должно производиться главным образом поперёк этого измерения, а не вдоль него. Какого рода замощения можно было бы предложить, чтобы извлечь пользу из этого априорного знания?

\begin{solution}
    Пусть важная ось (вдоль которой функция ценности меняется быстро) --- это
    $y$, а маловажная --- $x$. Предлагаются плитки в виде узких горизонтальных
    полос: вытянутых вдоль $x$ и узких по $y$.

    Такая форма сразу даёт нужную асимметрию обобщения. При смещении вдоль $x$
    (вдоль полосы) состояние почти всегда остаётся в той же плитке, набор
    активных признаков не меняется, поэтому оценка вдоль $x$ почти постоянна ---
    широкое обобщение по маловажной оси. При смещении по $y$ (поперёк полос)
    из-за их узости состояние быстро переходит из одной полосы в другую, активные
    плитки меняются, и оценка может меняться резко --- слабое обобщение, высокое
    разрешение по важной оси.

    Форма плитки задаёт обобщение, но итоговое разрешение определяется тем,
    насколько густо наслаиваются смещённые замощения. Поэтому замощения нужно
    сдвигать в основном вдоль $y$: чем гуще сдвиг по $y$, тем больше различимых
    комбинаций активных плиток вдоль этой оси и тем тоньше разрешение там, где оно
    важно. Вдоль $x$ сдвиги не нужны или редки --- там разрешение намеренно грубое,
    и лишние замощения только тратили бы память.

    Итого: узкие горизонтальные полосы, смещаемые преимущественно вдоль $y$ и
    почти не смещаемые вдоль $x$. Это даёт широкое обобщение по маловажной оси $x$
    и тонкое разрешение по важной оси $y$ --- ровно то, что подсказывает априорное
    знание.
\end{solution}

\subsection{Упражнение 9.5(с. 266)}
Предположим, что плиточное кодирование используется для преобразования
семимерного непрерывного пространства состояний в бинарные векторы признаков с
целью оценивания функции ценности состояния $\hat v(s,\mathbf w)\approx
v_\pi(s)$. Есть основания полагать, что зависимость между измерениями слабая,
поэтому мы решаем использовать восемь замощений по каждому измерению в
отдельности (полосами), т.\,е. всего $7\times 8 = 56$ замощений. Кроме того, на
случай, если между измерениями имеются какие-то попарные зависимости, мы берём
все $21$ пару измерений и для каждой пары строим по два замощения прямоугольными
плитками. Всего получается $21\times 2 + 56 = 98$ замощений. Имея эти векторы
признаков, мы всё же подозреваем, что следует применить усреднение для устранения
шума, поэтому решаем, что обучение должно быть постепенным, т.\,е. один и тот же
вектор признаков должен предъявляться примерно $10$ раз, прежде чем обучение
выйдет на асимптоту. Как следует выбрать размер шага $\alpha$? Объясните свой
ответ.

\begin{solution}
    Используем эвристику (9.19): $\alpha = \big(\tau\,\mathbb E[\mathbf x^\top
    \mathbf x]\big)^{-1}$. Здесь важно не число замощений само по себе, а
    $\mathbf x^\top\mathbf x$ --- число активных признаков (единиц в векторе).

    В каждом замощении, независимо от формы плиток (полоса или прямоугольник),
    состояние попадает ровно в одну плитку и активирует ровно один признак.
    Поэтому число активных признаков равно числу замощений:
    \begin{equation*}
        \mathbf x^\top\mathbf x = 56 + 21\cdot 2 = 98.
    \end{equation*}
    То, что прямоугольные замощения «двумерные», на это число не влияет.

    Причём $\mathbf x^\top\mathbf x$ здесь не случайно, а тождественно равно $98$ в
    любом состоянии (по одной активной плитке на замощение всегда). Это ровно тот
    идеальный случай из текста, когда $\mathbf x^\top\mathbf x$ постоянно, так что
    $\mathbb E[\mathbf x^\top\mathbf x] = 98$ и эвристика точна, а не приближённа.

    Подставляя $\tau = 10$:
    \begin{equation*}
        \alpha = \frac{1}{\tau\,\mathbf x^\top\mathbf x} = \frac{1}{10\cdot 98}
        = \frac{1}{980}.
    \end{equation*}
\end{solution}
